\chapter{Planificación y gestión del proyecto}\label{appendix:gestion}

Este apéndice documenta la metodología de trabajo, las herramientas empleadas y la planificación
temporal del desarrollo del sistema descrito en el Capítulo~\ref{chapter:desarrollo}.

\section{Metodología y ciclo de vida}\label{section:gestion_metodologia}

El proyecto siguió un enfoque ágil inspirado en \englishText{Scrum}, organizando el trabajo en siete
iteraciones funcionales alineadas con las etapas del \englishText{pipeline} de procesamiento clínico
(Tabla~\ref{tab:iteraciones_pipeline} del Capítulo~\ref{chapter:desarrollo}). Cada iteración
produjo un incremento verificable antes de abordar la siguiente etapa, reduciendo el riesgo de
integración y permitiendo ajustar la arquitectura de forma incremental.

\section{Herramientas y plataformas}\label{section:gestion_herramientas}

\begin{itemize}
    \item \textbf{Desarrollo:} Visual Studio Code / Cursor, Git.
    \item \textbf{Cliente móvil:} React Native, Expo, \texttt{expo-audio}.
    \item \textbf{\englishText{Gateway}:} Node.js, Express, \texttt{ws}.
    \item \textbf{Servicio \englishText{ASR}:} Python~3, FastAPI, Uvicorn, \englishText{faster-whisper},
    \englishText{Hugging Face Transformers}.
    \item \textbf{Terminología:} Docker Compose, Elasticsearch, Snowstorm.
    \item \textbf{Generación:} Ollama (\texttt{llama3.1:8b}).
    \item \textbf{Documentación:} \LaTeX, TikZ.
\end{itemize}

\section{Planificación. Diagrama de Gantt}\label{section:gestion_gantt}

La Tabla~\ref{tab:gantt_pipeline} presenta la secuenciación temporal de las iteraciones de
desarrollo. Las duraciones son orientativas y reflejan el orden de dependencias entre etapas:
cada iteración requiere la estabilización de la anterior.

\begin{table}[htbp]
    \centering
    \caption{Planificación temporal del desarrollo}
    \label{tab:gantt_pipeline}
    \small
    \begin{tabular}{@{}p{5.5cm}cccccccc@{}}
    \toprule
    \textbf{Tarea / Iteración} &
    \textbf{S.1} & \textbf{S.2} & \textbf{S.3} & \textbf{S.4} &
    \textbf{S.5} & \textbf{S.6} & \textbf{S.7} & \textbf{S.8} \\
    \midrule
    1. Captura y envío de audio (móvil + \englishText{gateway}) &
    \cellcolor{blue!15} & \cellcolor{blue!15} & & & & & & \\
    2. Transcripción \englishText{ASR} por fragmentos &
    & \cellcolor{blue!15} & \cellcolor{blue!15} & & & & & \\
    3. \englishText{NER} dual y fusión de entidades &
    & & \cellcolor{blue!15} & \cellcolor{blue!15} & & & & \\
    4. Enriquecimiento \englishText{SNOMED CT} y \englishText{CIE-10-ES} &
    & & & \cellcolor{blue!15} & \cellcolor{blue!15} & & & \\
    5. Reporte de control e informe \englishText{LLM} &
    & & & & \cellcolor{blue!15} & \cellcolor{blue!15} & & \\
    6. \englishText{Bundle FHIR} R4B &
    & & & & & \cellcolor{blue!15} & \cellcolor{blue!15} & \\
    7. Validación \englishText{end-to-end} y redacción del TFM &
    & & & & & & \cellcolor{blue!15} & \cellcolor{blue!15} \\
    \bottomrule
    \end{tabular}
    \end{table}
    Las columnas S.1--S.8 representan sprints de desarrollo consecutivos del periodo de desarrollo del trabajo fin de
    máster. Cada sprint tiene una duración de 2 semanas. Las celdas sombreadas indican el intervalo de ejecución de cada iteración.

\section{Sostenibilidad y mantenimiento}\label{section:gestion_sostenibilidad}


La arquitectura modular basada en el patron de diseño \englishText{Factory} y configuración externa (\texttt{config.yaml}) facilita
el mantenimiento evolutivo: la sustitución de modelos \englishText{ASR}, \englishText{PLN} o
\englishText{LLM} no requiere reescribir la capa de orquestación. Los costes de despliegue se
limitan al hardware local (CPU/GPU para inferencia, almacenamiento para Snowstorm) y al tiempo de
configuración de los servicios Docker y Ollama, sin dependencia de APIs de pago por uso para el
procesamiento de datos clínicos.
